\chapter{Umsetzung}
\label{ch:umsetzung}

% Vorgabe: Abschnitte 5.3-5.5 an die tatsaechlichen Kernmodule des eigenen
% Projekts anpassen (umbenennen, Anzahl variieren) -- nicht wortwoertlich
% uebernehmen. Interessante Quelltext-Ausschnitte hier per \lstinputlisting
% einbinden oder auf Anlage~\ref{anh:quellcode} verweisen; vollstaendiger
% Quelltext gehoert NICHT in den Hauptteil (Seitenlimit).

\section{Datenbank-/Datenmodelldesign}
[PLATZHALTER: Datenmodell, Tabellen-/Entitaetsstruktur. Vollstaendiges
ER-Modell/Tabellenmodell als Anlage~\ref{anh:diagramme} beigefuegt.]

\section{Projektinitialisierung und Grundkonfiguration}
[PLATZHALTER: Aufbau des Projektgeruests, Konfiguration, Basiskomponenten.]

\section{[PLATZHALTER: Kernmodul 1 -- Name einsetzen]}
[PLATZHALTER: Umsetzung des ersten fachlichen Kernbausteins.]

\section{[PLATZHALTER: Kernmodul 2 -- Name einsetzen]}
[PLATZHALTER: Umsetzung des zweiten fachlichen Kernbausteins.]

\section{Benutzeroberflaeche}
[PLATZHALTER: Umsetzung der UI/UX, Bezug zu den Mockups aus
Anlage~\ref{anh:mockups}. Screenshots der fertigen Anwendung als
Anlage~\ref{anh:screenshots} beigefuegt.]

\section{Hilfsfunktionen und weitere Logik}
[PLATZHALTER: Uebergreifende Hilfsfunktionen, Fehlerbehandlung, Logging o.Ae.]
